\chapter{Algebraic Generating Functions and the $n^{-3/2}$ Law}

Chapter~4 explained how a square-root singularity leads to an $n^{-3/2}$ coefficient law. This chapter studies one important source of such singularities: nonlinear algebraic systems coming from unambiguous context-free structure. The conclusion is striking, but it must be stated carefully. The exponent $-3/2$ is \emph{not} forced by being context-free all by itself. Simple unambiguous context-free languages can still have rational generating functions and pole-type asymptotics.

A good warning example is
\[
  a^* = \{\varepsilon, a, aa, aaa, \dots\},
\]
generated by the unambiguous grammar
\[
  S \to aS \mid \varepsilon.
\]
Its generating function is
\[
  S(z) = 1 + zS(z) = \frac{1}{1-z},
\]
so the coefficients are $1,1,1,\dots$, not of the form $Cn^{-3/2}\rho^{-n}$. Thus the square-root law belongs to a \emph{nonlinear generic regime}, not to every context-free language without exception.

The goal of the chapter is to explain the pipeline
\begin{align*}
  &\text{unambiguous grammar} \\
  &\Longrightarrow \text{algebraic generating function} \\
  &\Longrightarrow \text{generic square-root singularity} \\
  &\Longrightarrow n^{-3/2}\rho^{-n},
\end{align*}
while also making clear where the hidden hypotheses live.

\section{Unambiguous Context-Free Grammars}

A \emph{context-free grammar} is a tuple
\[
  \mathcal{G} = (V,\Sigma,S,\mathcal{P}),
\]
where $V$ is a finite set of nonterminals, $\Sigma$ is a finite terminal alphabet, $S \in V$ is the start symbol, and $\mathcal{P}$ is a finite set of productions of the form
\[
  A \to \alpha, \qquad A \in V, \quad \alpha \in (V \cup \Sigma)^*.
\]
The language $L(\mathcal{G})$ consists of all terminal strings derivable from $S$.

A grammar is \emph{unambiguous} if every word in the language has a unique parse tree. Equivalently --- a standard fact from formal-language theory --- every word has a unique leftmost derivation. This matters because our generating functions count \emph{distinct words}; if one word had several parse trees, a naive grammar-to-generating-function translation would overcount it.

\section{From Productions to Polynomial Systems}

The translation rule is straightforward once written symbol by symbol. Suppose a production has the form
\[
  A \to X_1 X_2 \cdots X_m,
\]
where each $X_i$ is either a terminal, a nonterminal, or the empty word $\varepsilon$.
\begin{itemize}
  \item A terminal contributes a factor of $z$ because it adds one unit of word length.
  \item A nonterminal $B$ contributes a factor of $B(z)$, the generating function for words derivable from $B$.
  \item The empty word contributes $1$.
\end{itemize}
Multiplying those contributions gives the monomial associated with the production. Summing over all productions with left-hand side $A$ gives an equation for $A(z)$.

This is exactly the symbolic-method logic from Chapter~3: concatenation becomes multiplication, alternative productions become addition, and unambiguity is what ensures the resulting count is by words rather than by derivations.

\begin{example}[$a^n b^n$]
The grammar
\[
  S \to aSb \mid \varepsilon
\]
generates $\{a^n b^n : n \ge 0\}$. The production $aSb$ contributes
\[
  z \cdot S(z) \cdot z = z^2 S(z),
\]
and the production $\varepsilon$ contributes $1$. So
\[
  S(z) = 1 + z^2 S(z),
  \qquad
  S(z) = \frac{1}{1-z^2}.
\]
This is rational, not square-root. Its coefficients are exactly $1$ for even $n$ and $0$ for odd $n$.
\end{example}

\begin{example}[Plane trees revisited]
The symbolic specification
\[
  \mathcal{T} = \mathcal{Z} \times \SEQ(\mathcal{T})
\]
from Chapter~3 is not literally a grammar production, but it obeys the same translation logic: a root contributes $z$ and the ordered list of children contributes $1/(1-T(z))$. Thus
\[
  T(z) = \frac{z}{1-T(z)},
\]
which rearranges to the polynomial equation
\[
  T(z)^2 - T(z) + z = 0.
\]
This is our first nonlinear algebraic example and will serve as the model case below.
\end{example}

After standard grammar normalization, one often obtains a finite polynomial system
\[
  Y_i(z) = \Phi_i(z,Y_1(z),\dots,Y_k(z)), \qquad i=1,\dots,k,
\]
with each $\Phi_i$ a polynomial having nonnegative coefficients. Chomsky normal form is a convenient way to keep the dependence at most quadratic in the nonterminal variables, though we will not need the full details of that conversion here.

\section{Why the Generating Function Is Algebraic}

The decisive structural theorem is the following.

\begin{theorem}[Chomsky--Sch\"utzenberger, used as a major black box]
If $L$ is generated by an unambiguous context-free grammar and $c_n$ denotes the number of words of length $n$ in $L$, then the ordinary generating function
\[
  L(z) = \sum_{n \ge 0} c_n z^n
\]
is algebraic over $\mathbb{Q}(z)$.
\end{theorem}

One standard proof translates the grammar into a polynomial system for the nonterminal generating functions and then uses elimination theory to extract a single polynomial equation for the start-symbol series. The elimination step is advanced algebra, so we treat it as a black box here. The key point to keep in mind is simpler: a finite unambiguous grammar creates a finite algebraic system, and that finiteness is what forces algebraicity.

There is also a uniqueness issue. Polynomial systems can have many formal solutions in principle, but the grammar singles out one nonnegative power-series solution by coefficient recursion in length: the coefficient of $z^n$ is determined by the number of derivations producing words of length $n$, and unambiguity prevents a hidden multiplicity problem.

\section{Puiseux Expansions: What Algebraic Singularities Look Like}

Algebraic functions do not have arbitrary singularities.

\begin{theorem}[Puiseux, used as a major black box]
Let $F(z)$ be an algebraic function satisfying a polynomial equation
\[
  P(z,F(z)) = 0.
\]
Near any singularity $\rho$ of a chosen branch of $F$, there is a positive integer $d$ and an expansion of the form
\[
  F(z) = \sum_{k \ge k_0} a_k \left(1-\frac{z}{\rho}\right)^{k/d},
\]
valid in a slit neighborhood of $\rho$.
\end{theorem}

What does this mean in practice?
\begin{itemize}
  \item The generating function defined by its power series at $0$ gives one distinguished branch.
  \item Near a singularity, that branch may continue into a local expression involving fractional powers.
  \item Integer exponents belong to the analytic part; the first genuinely fractional exponent is the first genuinely singular term.
\end{itemize}
Thus algebraic singularities are much more rigid than arbitrary singularities. They are poles or branch points of rational order, never essential singularities.

For quadratic equations, the generic fractional exponent is $1/2$, which is why square roots appear so often. But ``generic'' is the key word: algebraicity alone does not force a square root in every case, and certainly does not rule out rational degeneracies.

\section{Why Square Roots Are Generic in the Nonlinear Case}

We now explain the mechanism heuristically. This section is a local model, not a full proof.

Suppose a generating function $F(z)$ satisfies
\[
  P(z,F(z)) = 0,
\]
and suppose the relevant branch stays bounded up to its dominant singularity, with finite limit
\[
  \tau = \lim_{z \to \rho} F(z).
\]
If the implicit-function theorem were still active at $(\rho,\tau)$, then the branch would continue analytically past $\rho$. So a singularity must occur where
\[
  P(\rho,\tau)=0
  \qquad\text{and}\qquad
  P_y(\rho,\tau)=0.
\]
This pair of equations is called the \emph{characteristic system}.

Now write
\[
  \Delta z = z-\rho,
  \qquad
  \Delta F = F(z)-\tau.
\]
Expanding $P$ near $(\rho,\tau)$ gives, schematically,
\[
  0 = P_z\,\Delta z + P_y\,\Delta F + \frac{1}{2}P_{yy}(\Delta F)^2 + P_{zy}\,\Delta z\,\Delta F + \cdots.
\]
At the characteristic point, the linear term in $\Delta F$ vanishes because $P_y(\rho,\tau)=0$. If we now guess that the first nontrivial balance is
\[
  \Delta F = O(\sqrt{\Delta z}),
\]
then the mixed term $\Delta z\,\Delta F$ is of order $(\Delta z)^{3/2}$, smaller than both $\Delta z$ and $(\Delta F)^2$. So the leading balance becomes
\[
  P_z\,\Delta z + \frac{1}{2}P_{yy}(\Delta F)^2 \approx 0.
\]
Provided
\[
  P_z(\rho,\tau) \ne 0
  \qquad\text{and}\qquad
  P_{yy}(\rho,\tau) \ne 0,
\]
one obtains
\[
  \Delta F \asymp K \sqrt{\Delta z}
\]
for some nonzero constant $K$. Since
\[
  \Delta z = -\rho\left(1-\frac{z}{\rho}\right),
\]
this is the same as saying that the singular term is some branch-dependent constant times
\[
  \sqrt{1-z/\rho}.
\]
That is the local origin of the square-root law.

The word \emph{generic} means that higher degeneracies require extra coincidences. If one also had $P_{yy}(\rho,\tau)=0$, then a cubic or still higher balance would appear instead. Such cases exist, but they are not the default.

\section{The Generic $n^{-3/2}$ Law}

The theorem actually used in analytic combinatorics is not ``every context-free language has a square-root singularity.'' It is a more carefully scoped theorem about nonlinear positive algebraic systems.

\begin{theorem}[Generic square-root law for nonlinear context-free systems]
\label{thm:n32law}
Consider an unambiguous context-free grammar whose counting generating functions form a nonlinear polynomial system with nonnegative coefficients. Assume the standard structural hypotheses from the Drmota--Lalley--Woods type theory: the relevant dependency graph is strongly connected / irreducible, the system is proper and genuinely nonlinear, and the resulting coefficient sequence is aperiodic. Then the start-symbol generating function has a unique dominant singularity $\rho > 0$ of square-root type,
\[
  L(z) = g(z) - h(z)\sqrt{1-z/\rho},
\]
with $g$ and $h$ analytic near $\rho$ and $h(\rho) > 0$. Consequently,
\[
  c_n \sim C\,n^{-3/2}\rho^{-n}
\]
for some constant $C>0$.
\end{theorem}

Informally, these hypotheses mean the following. Build a dependency graph on the nonterminals, with an edge $A \to B$ whenever $B$ appears on the right-hand side of a production for $A$. ``Strongly connected'' or ``irreducible'' means every nonterminal can influence every other through this graph. ``Genuinely nonlinear'' means the polynomial system is not merely linear in the unknown generating functions; some right-hand side must contain a product of unknown series, so rational degeneracies like $a^*$ are excluded. ``Aperiodic'' means
\[
  \gcd\{n \ge 0 : c_n \ne 0\} = 1.
\]
Finally, ``proper'' is a shorthand for the standard well-foundedness assumptions excluding pathological epsilon-cycles and ensuring that the coefficient recursion is meaningful.

This theorem deliberately excludes the linear and rational cases illustrated earlier. Languages such as $a^*$ or $\{a^n b^n\}$ remain context-free and unambiguous, but their generating functions are rational, so their coefficient asymptotics are controlled by poles rather than square-root branch points.

\section{Applying the Transfer Theorem}

Assume we are in the setting of Theorem~\ref{thm:n32law}, so that
\[
  L(z) = g(z) - h(z)\sqrt{1-z/\rho}.
\]
The singular information sits in the second term. By Chapter~4,
\[
  [z^n]\left(1-z/\rho\right)^{1/2}
  \sim -\frac{1}{2\sqrt{\pi}} n^{-3/2}\rho^{-n}.
\]
Therefore
\[
  [z^n]\left(-h(\rho)\sqrt{1-z/\rho}\right)
  \sim \frac{h(\rho)}{2\sqrt{\pi}} n^{-3/2}\rho^{-n}.
\]
In the theorem setting, the analytic remainder extends farther than the singular term and therefore contributes exponentially smaller coefficients, so
\[
  c_n \sim \frac{h(\rho)}{2\sqrt{\pi}} n^{-3/2}\rho^{-n}.
\]
This is the promised $n^{-3/2}$ law.

\section{Worked Example: Plane Trees Revisited}

The plane-tree equation from Chapter~3 is
\[
  T(z)^2 - T(z) + z = 0.
\]
Here
\[
  P(z,F) = F^2 - F + z.
\]
The characteristic system is
\[
  P(\rho,\tau)=0,
  \qquad
  P_F(\rho,\tau)=2\tau-1=0.
\]
So $\tau = 1/2$, and substituting into $P=0$ gives
\[
  \frac{1}{4} - \frac{1}{2} + \rho = 0,
  \qquad
  \rho = \frac{1}{4}.
\]
The local characteristic-system picture predicts a square-root singularity. The exact quadratic formula confirms this:
\[
  T(z) = \frac{1-\sqrt{1-4z}}{2} = \frac{1}{2} - \frac{1}{2}\sqrt{1-z/\rho}.
\]
So here
\[
  h(\rho) = \frac{1}{2}.
\]
The generic law gives
\[
  [z^n]T(z)
  \sim \frac{1/2}{2\sqrt{\pi}} n^{-3/2}4^n
  = \frac{4^n}{4\sqrt{\pi}n^{3/2}}.
\]
Since $[z^n]T(z) = C_{n-1}$, this agrees with the classical Catalan asymptotic
\[
  C_{n-1} \sim \frac{4^{n-1}}{\sqrt{\pi}(n-1)^{3/2}} \sim \frac{4^n}{4\sqrt{\pi}n^{3/2}}.
\]

\section{What This Does and Does Not Let Us Infer}

The $n^{-3/2}$ law is mathematically strong, but it should not be used carelessly.

\begin{itemize}
  \item If one proves that a counting generating function belongs to the nonlinear square-root regime above, then the exponent $-3/2$ is a theorem.
  \item If one sees a finite-data fit with exponent near $-3/2$, that is only \emph{heuristic evidence} for such a regime, not a proof.
  \item Other exponents point to other singular models. For instance, a simple pole corresponds to exponent $0$, a logarithmic singularity gives exponent $-1$, and an inverse square root gives exponent $-1/2$.
\end{itemize}

So empirical exponent fitting can be suggestive, but only once one also remembers the limitations of finite data, periodicity, and model mismatch. Chapter~19 will return to these issues in a more explicit research-roadmap style.

\begin{remark}
Aperiodicity is not the only hypothesis that matters. Linear / rational cases are another important source of exceptions to the square-root law. So ``aperiodic'' should never be read as a synonym for ``therefore $n^{-3/2}$.''
\end{remark}

\section*{Summary of Part I}

Chapters~1--5 together provide the basic analytic-combinatorics pipeline.
\begin{enumerate}
  \item Encode a combinatorial class by a generating function.
  \item Use structural equations to identify the relevant analytic object.
  \item Locate and classify the dominant singularity.
  \item Apply a transfer theorem to recover coefficient asymptotics.
\end{enumerate}

For a broad and important family of \emph{nonlinear} algebraic systems, the outcome is the universal-looking law
\[
  c_n \sim C n^{-3/2}\rho^{-n}.
\]
But this is a theorem about a specific regime, not about all context-free languages indiscriminately.

The rest of the book moves from pure counting toward weighted and probabilistic models. That shift will require new care, because counting OGFs, probability generating functions, and partition sums are related objects, not identical ones.

\section*{Further Reading}

\begin{itemize}
  \item \textbf{N. Chomsky and M.\,P. Sch\"utzenberger}, ``The algebraic theory of context-free languages,'' in \emph{Computer Programming and Formal Systems} (North-Holland, 1963) \cite{ChomskySchutzenberger1963} --- The original paper proving algebraicity for unambiguous context-free languages.

  \item \textbf{M. Drmota}, \emph{Random Trees} (Springer, Vienna, 2009) \cite{Drmota2009} --- A detailed source for characteristic systems, square-root singularities, and nonlinear algebraic systems.

  \item \textbf{C. Banderier and M. Drmota}, ``Formulae and asymptotics for coefficients of algebraic functions,'' \emph{Combinatorics, Probability and Computing} (2015) \cite{BanderierDrmota2015} --- Develops generic asymptotic laws for algebraic systems well beyond the simplest tree examples.

  \item \textbf{P. Flajolet and R. Sedgewick}, \emph{Analytic Combinatorics} (Cambridge University Press, 2009) \cite{FlajoletSedgewick2009}, Chapter~VII --- Covers algebraic functions, Newton polygons, and Puiseux expansions from the analytic-combinatorics point of view.
\end{itemize}
